% 09-use.tex: holders, relying parties and the broker.
\section{Using the credential: holders, relying parties and the login that keeps no record}
\label{sec:use}

\subsection{The wallet is a protocol, not a product}

The reference holder agent, \code{scripts/polaris-wallet.py}, holds a credential as a file, shows
it, verifies it offline through the detached verifier, presents it, proves membership in zero
knowledge against a published epoch, signs a document on the holder's behalf and logs in to a
relying party. Every message it emits is specified so that any client can emit the same. Native
mobile and desktop clients, card middleware and a browser bridge are product engineering that a
reference implementation on notional data does not attempt; they would follow the protocol, and
never lead it. \sImpl{} for the protocol and the reference script; \sFut{} for native clients.

A presentation is an unsigned wrapper around signed things: the issuer-signed credential, an
optional stapled status assertion so that authorization is decidable with no connectivity, an
optional membership proof, the context and disclosure level presented under, and an opaque
presentation code. The verifier decides it offline and reports the code as present or absent
without ever interpreting it, because a duress code is matched only by the issuer, out of sight; a
duress presentation is byte-for-byte the same shape as a normal one. For transfer by QR or NFC the
presentation is compressed and split into frames that each name the digest of the whole; a receiver
takes them in any order and refuses, before parsing anything, frames naming different digests, a
missing index, a conflicting duplicate, or a reassembled payload whose digest does not match. The
decoder is resource-bounded: at most 9,999 frames, 512~KiB compressed and 2~MiB inflated, so a
decompression bomb is refused at the limit rather than decompressed. The bounds were added after an
outside review named the gap. \sImpl

\subsection{The relying party}

The reference relying party (\code{scripts/polaris-relying-party.py}) decides accept or reject for a
presentation by combining offline authenticity with online status, refusing a revoked, tampered or
foreign-issuer credential and staying blind to duress, and an end-to-end drill runs the whole
wallet-to-relying-party matrix under real signatures every release. A relying party registers with
an issuer as an organisation with a scope that a schema constraint bounds to exactly \code{verify},
\code{authenticate}, or both, so that a fourth kind of access is a constitutional change rather than
a configuration one. Identity never becomes a login product by accident. \sImpl

\subsection{Logging in without a login record}

The auth broker lets a relying party offer \emph{log in with a Polaris credential} through an
authorization code with PKCE. The holder presents the credential at its issuing authority's
instance by possession, and the relying party exchanges the resulting code for an identity token
signed by the issuing agency, which it verifies offline. Four guards keep this from becoming a login
record. The token's subject is the credential's hash, never a person identifier, and it is stable
per credential and therefore correlatable across relying parties, which the repository documents as
a permanent property rather than pretending otherwise. The token carries no claim beyond the
vocabulary: context, disclosure level, the assurance reached, enrolment status, instants. The
authorize route writes nothing, and the token endpoint's only write is the code's hash into a
register that holds no subject and no relying party. And duress is served identically. \sImpl

An outside review at \code{v9.331} found that the relying party's demands, a required
zero-knowledge step-up or a required enrolment status, arrived only in the holder-side request, so
nothing bound the authority to what the relying party actually required. Since \code{v9.336} the
policy is registered on the relying party's row and applied first: a request may add a requirement
and never remove one, and a context other than the registered one is refused. The same review noted
that the authorization code was signed but readable; it is now encrypted, because before this flow
is ever used through a browser rather than the wallet, an opaque code is the floor. \sImpl

\subsection{What is deliberately not built}

There is no session at the authority, no consent screen, no discovery document, and no browser
redirect choreography: a relying party integrates from the wire specification and the registry. A
WebAuthn browser bridge waits on a holder-side key that the issuer-centric model does not have; the
holder holds a credential, not a key pair, which is why document signing is notarial and login is
by possession. No relying party outside the repository has integrated. \sExt

\begin{layercard}{Layer card: use}
\qa{What it is}{The wallet protocol and reference script, presentations and QR framing, the reference relying party, relying-party registration with a bounded scope, and the auth broker.}
\qa{Why it exists}{A credential nobody can hold, present or accept is a database row. This layer is where the credential meets software the issuer does not run.}
\qa{What it trusts}{The signed objects inside a presentation; the issuer's key for the identity token; the relying party's registered policy.}
\qa{What it refuses}{To interpret the presentation code; to write a login record; to accept a scope outside the CHECK; to let a holder's request weaken the relying party's registered demand; to decompress an unbounded payload.}
\qa{What it can see}{A relying party sees the credential value and the verdicts. The authority sees a possession proof and consumes a code hash.}
\qa{What it cannot see}{The authority cannot see who logged in where; the relying party cannot see the person's record.}
\qa{If it fails}{A replayed code hits the consumed register; a wrong presentation, an unmet step-up, a verify-only bearer at the broker, a stale or tampered frame are each refused in the drills.}
\qa{Checked by}{\code{check\_holder\_wallet}, \code{check\_wallet\_presentation}, \code{check\_holder\_verifier\_flow}, \code{check\_auth\_broker}, \code{check\_relying\_party\_api}, \code{check\_qr\_resource\_bounds}.}
\qa{Now or later}{Now, as a protocol. Native clients and a real relying party are later.}
\qa{Inside and outside}{Inside: the privacy boundary. Outside: federation, because a relying party rarely trusts only one authority.}
\end{layercard}

\nextlayer{A relying party that trusts one authority is easy. A society has many authorities, none
of which should be able to speak for the others, and none of which should hold every person's
record. The next layer is how independent authorities recognise one another without collapsing into
one.}
